% Source PDF page 15 (printed p. 1)
本章概述数学优化，重点介绍凸优化的特殊地位。这里非正式引入的概念，将在后续章节中以更严谨、更详细的方式讨论。

\section{数学优化}
\emph{数学优化问题}，简称\emph{优化问题}，具有如下形式：
\begin{equation}
\begin{array}{ll}
\text{最小化} & f_0(x)\\
\text{满足} & f_i(x)\le b_i,\quad i=1,\ldots,m.
\end{array}\tag{1.1}
\end{equation}
这里，向量 $x=(x_1,\ldots,x_n)$ 是问题的\emph{优化变量}，函数 $f_0:\R^n\to\R$ 是\emph{目标函数}，函数 $f_i:\R^n\to\R$（$i=1,\ldots,m$）是（不等式）\emph{约束函数}，常数 $b_1,\ldots,b_m$ 是各约束的限制值或界。如果向量 $x^\star$ 在所有满足约束的向量中具有最小的目标值，就称它是\emph{最优的}，或称为问题（1.1）的一个\emph{解}。也就是说，对任意满足 $f_1(z)\le b_1,\ldots,f_m(z)\le b_m$ 的 $z$，都有 $f_0(z)\ge f_0(x^\star)$。

我们通常考察由目标函数和约束函数的特定形式所确定的优化问题族或问题类。一个重要的例子是\emph{线性规划}：如果优化问题（1.1）的目标函数与约束函数 $f_0,\ldots,f_m$ 都是线性的，即对所有 $x,y\in\R^n$ 和所有 $\alpha,\beta\in\R$，满足
\begin{equation}
 f_i(\alpha x+\beta y)=\alpha f_i(x)+\beta f_i(y),\tag{1.2}
\end{equation}
就称其为线性规划。如果优化问题不是线性的，则称为\emph{非线性规划}。

本书研究一类称为\emph{凸优化问题}的优化问题。凸优化问题的目标函数和约束函数都是凸函数，也就是说，它们满足不等式
\begin{equation}
 f_i(\alpha x+\beta y)\le\alpha f_i(x)+\beta f_i(y),\tag{1.3}
\end{equation}
其中 $x,y\in\R^n$ 任意，而 $\alpha,\beta\in\R$ 满足 $\alpha+\beta=1$、$\alpha\ge0$、$\beta\ge0$。

% Source PDF page 16 (printed p. 2)
比较（1.3）和（1.2）可以看出，凸性比线性更一般：不等式代替了限制更强的等式，而且只要求对某些 $\alpha$ 和 $\beta$ 成立。因此，每个线性规划都是凸优化问题，凸优化可以看作线性规划的推广。

\subsection{应用}
优化问题（1.1）抽象描述了这样一个问题：从一组候选选择中选出一个最好的 $\R^n$ 中的向量。变量 $x$ 表示作出的选择；约束 $f_i(x)\le b_i$ 表示限制可选方案的刚性要求或技术规范；目标值 $f_0(x)$ 表示选择 $x$ 的代价。（也可以将 $-f_0(x)$ 理解为选择 $x$ 的价值或效用。）优化问题（1.1）的解，就是在所有满足刚性要求的选择中代价最小（或效用最大）的那个选择。

例如，在投资组合优化中，我们希望找到将一笔资金投向 $n$ 项资产的最佳方式。变量 $x_i$ 表示对第 $i$ 项资产的投资，因此向量 $x\in\R^n$ 描述了整个投资组合在各项资产间的配置。约束可以表示预算限制（即投资总额的限制）、投资非负的要求（假设不允许卖空），以及整个投资组合的最低可接受预期收益。目标函数或代价函数可以用来度量组合收益的整体风险或方差。在这种情况下，优化问题（1.1）就是在所有满足刚性要求的配置中，选择风险最小的投资组合配置。

另一个例子是电子设计中的器件尺寸设计，即选择电路中每个器件的宽度和长度。这里的变量表示器件的宽度和长度。约束表示各种工程要求，例如制造工艺对器件尺寸的限制、保证电路能够以指定速度可靠工作的时序要求，以及电路总面积的限制。器件尺寸设计问题中，一个常用的目标是电路消耗的总功率。优化问题（1.1）就是寻找满足设计要求（可制造性、时序和面积要求）且功率效率最高的器件尺寸。

在数据拟合中，任务是从一族候选模型中，找出与观测数据及先验信息最相符的模型。这里的变量是模型参数，约束可以表示先验信息或对参数的必要限制（例如非负性）。目标函数可以度量观测数据与模型预测值之间的不匹配程度或预测误差，也可以从统计上度量参数取值不太可能或不太合理的程度。优化问题（1.1）就是找出与先验信息一致、且对观测数据给出最小不匹配程度或预测误差的模型参数值（或者在统计框架下最有可能的参数值）。

% Source PDF page 17 (printed p. 3)
种类繁多的实际决策问题（以及系统设计、分析和运行问题），都可以写成数学优化问题，或者多准则优化问题等变体。事实上，数学优化已经成为许多领域的重要工具。它广泛应用于工程中的电子设计自动化、自动控制系统，以及土木、化工、机械和航空航天工程中的最优设计问题。优化还用于网络设计与运行、金融、供应链管理、调度等众多领域，应用范围仍在不断扩展。

在这些应用中，大多数时候数学优化是人类决策者、系统设计者或系统操作者的辅助工具。人负责监督过程、检查结果，并在必要时修改问题或求解方法。优化问题所建议的行动也由人来执行，例如买卖资产以实现最优投资组合。

一种较新的现象，为数学优化开辟了更多应用的可能。随着产品中嵌入式计算机日益普及，\emph{嵌入式优化}迅速发展。在这类应用中，优化用于自动作出实时选择，甚至执行相应操作，完全不需要或只需要很少的人为干预或监督。在一些领域，传统自动控制系统与嵌入式优化的融合已经相当深入；在另一些领域，这一过程才刚刚开始。嵌入式实时优化带来了一些新的挑战，尤其要求求解方法极其可靠，并能在可预测的时间内、使用可预测的存储量求解问题。

\subsection{求解优化问题}
一类优化问题的\emph{求解方法}，是指这样一种算法：给定这类问题中的某个具体问题，也就是一个\emph{问题实例}，算法便计算出达到给定精度的解。自20世纪40年代末以来，人们为各种优化问题开发算法、分析算法性质并实现高质量软件，投入了大量工作。这些算法的有效性，也就是求解优化问题（1.1）的能力，差别很大，取决于许多因素，例如目标函数和约束函数的具体形式、变量和约束的数量，以及稀疏性等特殊结构。（若每个约束函数只依赖少数几个变量，就称问题是稀疏的。）

即使目标函数和约束函数光滑（例如均为多项式），一般优化问题（1.1）仍然出乎意料地难以求解。因此，处理一般问题的方法都需要某种妥协，例如接受很长的计算时间，或者接受可能找不到解。第1.4节将讨论其中一些方法。

不过，“大多数优化问题很难求解”这一普遍规律也有一些重要的例外。对于少数几类问题，我们有高效算法，能够可靠地求解包含数百乃至数千个变量和约束的大规模问题。

% Source PDF page 18 (printed p. 4)
两个重要且广为人知的例子是最小二乘问题和线性规划，下面第1.2节将作介绍，第4章再详细讨论。凸优化也是这一规律的例外，只是较少为人所知：与最小二乘和线性规划一样，凸优化拥有十分高效的算法，能够可靠而高效地求解大规模问题。

\section{最小二乘与线性规划}
本节介绍凸优化中两个应用广泛、广为人知的特殊子类：最小二乘和线性规划。（第4章将对它们作完整、严格的讨论。）

\subsection{最小二乘问题}
\emph{最小二乘问题}是没有约束（即 $m=0$），且目标为若干形如 $a_i^{\T}x-b_i$ 的项的平方和的优化问题：
\begin{equation}
 \text{最小化}\quad f_0(x)=\norm{Ax-b}_2^2=\sum_{i=1}^k(a_i^{\T}x-b_i)^2.\tag{1.4}
\end{equation}
这里 $A\in\R^{k\times n}$（$k\ge n$），$a_i^{\T}$ 是 $A$ 的各行，向量 $x\in\R^n$ 是优化变量。

\subsubsection*{求解最小二乘问题}
最小二乘问题（1.4）的求解可以归结为求解一组线性方程
\[
 (A^{\T}A)x=A^{\T}b,
\]
因此有解析解 $x=(A^{\T}A)^{-1}A^{\T}b$。对于最小二乘问题，我们已有良好的算法及其软件实现，能够以极高的可靠性求得高精度解。求解时间近似正比于 $n^2k$，其中的比例常数是已知的。一台当时的台式计算机，几秒钟就能求解包含数百个变量和数千个平方项的最小二乘问题；更强的计算机当然能够求解更大的问题，或者更快地求解同样规模的问题。（此外，根据 Moore 定律，这些求解时间在未来还会按指数速度缩短。）最小二乘算法及其软件的可靠性，足以支持嵌入式优化。

许多情况下，利用系数矩阵 $A$ 的某种特殊结构，我们还能求解更大的最小二乘问题。例如，假设 $A$ 稀疏，也就是非零元素远少于 $kn$ 个。利用稀疏性，通常可以把求解时间降到远小于 $n^2k$ 的量级。一台当时的台式计算机，大约一分钟就能求解包含数万个变量和数十万个平方项的稀疏最小二乘问题，尽管具体时间取决于稀疏模式。

% Source PDF page 19 (printed p. 5)
对于极大规模的问题（例如包含数百万个变量），或者实时计算要求极其严格的问题，求解最小二乘问题可能仍有挑战。但在绝大多数情况下，现有方法非常高效，也极其可靠。可以说，对于尚未接近当前计算能力极限的最小二乘问题，其求解已经是一项成熟技术，许多不了解、也无须了解具体细节的人都能可靠地使用它。

\subsubsection*{使用最小二乘}
最小二乘问题是回归分析、最优控制以及许多参数估计和数据拟合方法的基础。它有多种统计解释，例如在线性观测受到高斯测量误差污染时，对向量 $x$ 作最大似然估计。

识别一个优化问题是否为最小二乘问题很直接：只需确认目标是二次函数，再检验相应的二次型是否半正定。尽管基本最小二乘问题的形式简单固定，仍有一些标准技巧可以提高它在应用中的灵活性。

在\emph{加权最小二乘}中，我们最小化加权最小二乘代价
\[
 \sum_{i=1}^k w_i(a_i^{\T}x-b_i)^2,
\]
其中 $w_1,\ldots,w_k$ 为正。（这个问题很容易写成标准最小二乘问题并求解。）权重 $w_i$ 用来表示对各项 $a_i^{\T}x-b_i$ 大小的关注程度不同，也可以仅仅用来影响解。在统计中，加权最小二乘出现在由方差不等的误差污染的线性观测估计向量 $x$ 的问题中。

最小二乘中的另一种技巧是\emph{正则化}，即在代价函数中增加额外项。最简单的情形是在代价中加入变量平方和的一个正倍数：
\[
 \sum_{i=1}^k(a_i^{\T}x-b_i)^2+\rho\sum_{i=1}^n x_i^2,
\]
其中 $\rho>0$。（这个问题同样可以写成标准最小二乘问题。）这些额外项惩罚 $x$ 的较大取值；在只最小化第一项无法得到合理解的情况下，它们可以给出合理的解。参数 $\rho$ 由使用者选择，用来权衡原目标函数 $\sum_{i=1}^k(a_i^{\T}x-b_i)^2$ 足够小与 $\sum_{i=1}^n x_i^2$ 不太大这两个要求。在统计估计中，若给待估向量 $x$ 指定一个先验分布，便会出现正则化。

第6章介绍加权最小二乘和正则化，第7章给出它们的统计解释。

% Source PDF page 20 (printed p. 6)
\subsection{线性规划}
另一类重要的优化问题是\emph{线性规划}，其目标函数和所有约束函数都是线性的：
\begin{equation}
\begin{array}{ll}
\text{最小化}&c^{\T}x\\
\text{满足}&a_i^{\T}x\le b_i,\quad i=1,\ldots,m.
\end{array}\tag{1.5}
\end{equation}
这里向量 $c,a_1,\ldots,a_m\in\R^n$ 和标量 $b_1,\ldots,b_m\in\R$ 是确定目标函数及约束函数的问题参数。

\subsubsection*{求解线性规划}
线性规划的解没有像最小二乘那样简单的解析公式，但有多种非常有效的求解方法，包括 Dantzig 的单纯形法，以及本书后面将介绍的较新的内点法。尽管无法像最小二乘那样给出求解所需算术运算次数的精确值，我们仍能对内点法达到给定精度所需的运算次数建立严格上界。实际复杂度为 $n^2m$ 阶（假设 $m\ge n$），只是其中常数不像最小二乘那样明确。这些算法相当可靠，虽然可能略逊于最小二乘方法。使用小型台式计算机，几秒钟就能轻松求解包含数百个变量和数千个约束的问题。如果问题稀疏，或者有其他可利用的结构，通常还能求解包含数万乃至数十万个变量和约束的问题。

与最小二乘一样，极大规模线性规划，或者实时计算要求极其严格的线性规划，仍有求解上的挑战。不过，求解大多数线性规划也可以说已经是一项成熟技术。线性规划求解器能够嵌入、也确实已嵌入许多工具和应用中。

\subsubsection*{使用线性规划}
有些应用直接得到形式为（1.5）或其他若干标准形式之一的线性规划。还有许多情况下，原优化问题并非标准线性规划形式，但可以用第4章详细介绍的技巧转化为等价的线性规划，然后求解。

一个简单的例子是 Chebyshev 逼近问题：
\begin{equation}
 \text{最小化}\quad\max_{i=1,\ldots,k}\abs{a_i^{\T}x-b_i}.\tag{1.6}
\end{equation}
这里 $x\in\R^n$ 是变量，$a_1,\ldots,a_k\in\R^n$ 和 $b_1,\ldots,b_k\in\R$ 是确定问题实例的参数。它与最小二乘问题（1.4）很相似：两者的目标都度量各项 $a_i^{\T}x-b_i$ 的大小。最小二乘使用这些项的平方和，而 Chebyshev 逼近使用绝对值的最大值。

% Source PDF page 21 (printed p. 7)
另一个重要区别是，Chebyshev 逼近问题（1.6）的目标函数不可微，而最小二乘问题（1.4）的目标是二次函数，因此可微。

Chebyshev 逼近问题（1.6）可以通过求解下面的线性规划来求解：
\begin{equation}
\begin{array}{ll}
\text{最小化}&t\\
\text{满足}&a_i^{\T}x-t\le b_i,\quad i=1,\ldots,k,\\
&-a_i^{\T}x-t\le-b_i,\quad i=1,\ldots,k,
\end{array}\tag{1.7}
\end{equation}
变量为 $x\in\R^n$ 和 $t\in\R$。（细节见第6章。）线性规划容易求解，因此 Chebyshev 逼近问题也容易求解。

任何熟悉线性规划的人，都会看出 Chebyshev 逼近问题（1.6）可以归结为线性规划。但缺少这一背景的读者，可能不容易看出目标不可微的 Chebyshev 逼近问题（1.6）也能写成线性规划并求解。

识别可以归结为线性规划的问题，比识别最小二乘问题稍复杂，但只需掌握几个标准技巧，因此很容易学会。这项工作甚至可以部分自动化：有些用于描述和求解优化问题的软件系统，能够自动识别某些可重写为线性规划的问题。

\section{凸优化}
\emph{凸优化问题}具有如下形式：
\begin{equation}
\begin{array}{ll}
\text{最小化}&f_0(x)\\
\text{满足}&f_i(x)\le b_i,\quad i=1,\ldots,m,
\end{array}\tag{1.8}
\end{equation}
其中函数 $f_0,\ldots,f_m:\R^n\to\R$ 是凸函数，即对所有 $x,y\in\R^n$，以及满足 $\alpha+\beta=1$、$\alpha\ge0$、$\beta\ge0$ 的所有 $\alpha,\beta\in\R$，都有
\[
 f_i(\alpha x+\beta y)\le\alpha f_i(x)+\beta f_i(y).
\]
最小二乘问题（1.4）和线性规划问题（1.5）都是一般凸优化问题（1.8）的特例。

\subsection{求解凸优化问题}
凸优化问题的解通常没有解析公式，但与线性规划一样，存在非常有效的求解方法。内点法在实际中效果很好，并且某些情况下可以证明：达到指定精度所需的运算次数，不超过问题维数的某个多项式。（第11章讨论这一点。）

% Source PDF page 22 (printed p. 8)
我们将看到，内点法求解问题（1.8）所需的步骤或迭代次数，几乎总在10到100之间。不利用问题中的稀疏性等结构时，每一步需要的运算次数量级为
\[
 \max\{n^3,n^2m,F\},
\]
其中 $F$ 是计算目标函数和约束函数 $f_0,\ldots,f_m$ 的一阶与二阶导数所需的代价。

与线性规划方法一样，这些内点法相当可靠。在当时的台式计算机上，至多几十秒就能轻松求解包含数百个变量和数千个约束的问题。利用稀疏性等问题结构，还能求解包含成千上万个变量和约束的更大问题。

目前还不能说，一般凸优化问题的求解已经像最小二乘或线性规划那样成为成熟技术。一般非线性凸优化的内点法研究仍十分活跃，何种方法最好还没有形成共识。但有理由期待，在未来几年内，求解一般凸优化问题也将成为一项技术。对于某些凸优化子类，例如第4章详细研究的二阶锥规划和几何规划，可以说其内点法正在接近成熟技术的阶段。

\subsection{使用凸优化}
至少从概念上说，使用凸优化与使用最小二乘或线性规划十分相似。如果能将问题写成凸优化问题，就能像求解最小二乘那样高效地求解它。稍加夸张地说，将一个实际问题写成凸优化问题，就已经解决了原问题。

不过，两者也有一些重要区别。最小二乘问题很容易识别，凸函数却可能难以识别。此外，变换凸问题的技巧远比变换线性规划的技巧多。因此，识别凸优化问题，或者识别能转化为凸优化的问题，可能相当有挑战。本书的主要目标，就是为读者提供完成这项工作所需的背景。掌握识别或建立凸优化模型的技能后，你会发现，令人意外地多的问题都能通过凸优化求解。

使用凸优化的挑战和技巧，主要在于识别问题和建立模型。一旦完成建模，求解就像最小二乘和线性规划那样，几乎成为一项现成技术。
